\documentclass[book.tex]{subfiles}

\begin{document}

\chapter{Maxwell's Equations}
\label{chap:maxwells_equations}
\noindent\rule{11cm}{0.4pt} \\
\textbf{Prerequisites:} \autoref{chap:electrostatics}, \autoref{chap:electric_current}  \\
\textbf{Difficulty Level:} ** \\
\noindent\rule{11cm}{0.4pt}
\vspace{5mm}

Maxwell's equations are a collection of coupled partial differential equations that relate the electric and magnetic fields to one another.

Let the electric and magnetic fields be given by
\begin{align}
E: &\mathbb{R}^3 \to T \mathbb{R}^3 \\
B: &\mathbb{R}^3 \to T \mathbb{R}^3
\end{align} 
Let $\rho: \mathbb{R}^3 \to \mathbb{R}$ be a charge density function which tracks the continuous amount of charge at a given location. 

Maxwell's laws are defined by four equations that we discuss in the following sections. The four equations are coupled to one another and jointly define the electrical and magnetic fields as a set of 4 coupled PDEs.

\section{Gauss's Law}

Gauss's law dictates how a given charge distribution induces an electric field. We state the law below in its differential form.

\begin{align}
\nabla \cdot E = \frac{\rho}{\epsilon_0}
\end{align} 

Gauss's law can be viewed as a statement that the electric charge contained within a region of space induces a "flow" outwards of the electric field as measured by the divergence of the electric field. 

If the electric field is given by a potential, Gauss's law turns into Poisson's equation
\begin{align}
\nabla^2 V &= - \frac{\rho}{\varepsilon_0}
\end{align}
    
\section{Gauss's Law for Magnetism}

Gauss's law for magnetism stands in contrast to Gauss's law. There is no concept of a "magnetic charge" that induces an outward flow in the magnetic field. In fact, the magnetic field must always have divergence 0.
\begin{align}
\nabla \cdot B = 0
\end{align}
\section{Faraday's Law of Induction}

The two versions of Gauss's law have treated the electric and magnetic fields in isolation. Faraday's law describes how a time varying magnetic field induces an electric field.
\begin{align}
\nabla \times E = - \frac{\partial B}{\partial t}
\end{align}
\section{Amp\`{e}re's Law}
Amp\`{e}re's law describes how both electric current and a time varying electric field induce a magnetic field
\begin{align}
\nabla \times B = \mu_0 \left (J + \epsilon_0 \frac{\partial E}{\partial t} \right )
\end{align}

\section{A Physical Theory for Maxwell's Equations}

In this section, we transform Maxwell's laws into a physical theory in our formulation. The state space is given by a tuple for the electric and magnetic fields
\begin{align}
\mathcal{S} &= (\mathbb{R}^3 \to T\mathbb{R}^3)^2 \\
(E, B) &\in \mathcal{S}
\end{align} 
The state is the pair of the current electric and magnetic fields $(E, B)$. We define the electromagnetic energy density as
\begin{align}
U = \frac{\epsilon_0 \|E\|^2}{2} + \frac{\|B\|^2}{2 \mu_0}
\end{align} 
and the energy flux $u$ as
\begin{align}
u &= \frac{E \times B}{\mu_0}
\end{align}
The energy density obeys the conservation equation
\begin{align}
\frac{\partial U}{\partial t} + \nabla \cdot u &= - E \cdot J
\end{align}

The evolution of the system is governed by the solution to Maxwell's equations. The parameters of the theory are given by
\begin{align}
\mathcal{W} &= \{E_0, B_0, \epsilon_0, \mu_0\}
\end{align}
where $E_0$ and $B_0$ are the initial electric and magnetic fields. The domain of applicability $\mathcal{C}$ is governed by the fact that the magnitude of the fields shouldn't get too large 

\begin{align}
|E_0|, |B_0| < C
\end{align}

The hyperparameters are given by \begin{align}
\mathcal{HP} = \{C, \epsilon \}
\end{align}

The observables $\mathcal{B}$ are defined as follows. Using a test charge we can measure the electric and magnetic field strengths at a point up to some noise 
\begin{align}
f(E, B, x) = (E(x) + \textrm{Normal}(0, \epsilon), B(x) + \textrm{Normal}(0, \epsilon))
\end{align}

\section{A Lagrangian for Maxwell's Equations}
Let's now introduce a Lagrangian view of this system. Let's start with a few definitions. Let $E$ be an electric field. Let $B$ be a magnetic field. The magnetic vector potential is defined as the solution to the following equation
\begin{align}
    B &= \nabla \times A
\end{align}
The scalar potential $\phi(\vec{r}, t)$ is then defined as the solution to the following equation.
\begin{align}
    E(\vec{r}, t)  &= - \nabla \phi - \frac{\partial A}{\partial t}
\end{align}

The minimal coupling Lagrangian of a charged particle with mass $m$ and charge $q$ is defined as
\begin{align}
    L = \frac{1}{2} m \vec{v}^T\vec{v} + q \vec{v}^TA - q \phi 
\end{align}
where $\vec{v} = \frac{d \vec{r}}{dt}$ is the velocity 

%\textcolor{red}{TODO: Add a definition of a physical theory that uses this lagrangian.}

\end{document}